\chapter{Theory of integration}

\paragraph{On terminology:} Antiderivatives and integrals are two different things that shouldn't be confused. Although, they will be connected with the fundamental theorem of calculus at the end of this chapter, which is often considered as the milestone of mathematical analysis.

Originally, derivatives and integrals were two completely different things. The ancient Greeks invented the integrals as a way to calculate areas bounded a non-rectangular perimeter. Then, Newton and Liebniz came along and invented derivatives as a measure of rate of change. There were no connections between derivatives and integrals until seventy years later when James Gregory discovered the fundamental theorem of calculus, which links derivatives and integrals together.

In \cref{sec:integrals}, I've glazed over the fundamental theorem of calculus like it is something that is \emph{obviously true}. I've treated it like something that is just defined that way, when really it's not. So, in this chapter, I'd like to take you through the theory of integration. Defining integrals as a measure of area under the curve, unrelated to derivatives, then link it with the derivative later using the fundamental theorem of calculus. Then finally, we shall discuss why \emph{antiderivatives} can be represented as integrals.

\section{Riemann integrals}

A partition \(\mathcal{P}\) of a closed subset, \({a, b}\) is an arranged list of number, \((x_0, x_1, \dots, x_n)\) such that
\begin{equation}
  x_0 = a, x_n = b \mathand x_0 < x_1 < \dots < x_n.
\end{equation}

\section{Measures}
